% !TeX root = ../technical_report.tex
\chapter{FastAPI REST Service and Server-Rendered HTMX Operator Dashboard}
\label{ch:fastapi_service_architecture}

% CHAPTER BRIEF FOR JULES / TECHNICAL WRITER:
% -------------------------------------------------------------
% Focus: Dual-mode web application architecture, REST endpoints, and HTMX interactive UI.
% Repository Files:
%   - src/taxon_vision/service/api.py
%   - src/taxon_vision/service/schemas.py
%   - src/taxon_vision/service/inference.py
%   - src/taxon_vision/service/routes/predict.py
%   - src/taxon_vision/service/routes/health.py
%   - src/taxon_vision/service/routes/train.py
%   - src/taxon_vision/service/routes/explain.py
%   - src/taxon_vision/service/routes/views.py
%   - src/taxon_vision/service/templates/dashboard.html
% Detailed Technical Concepts:
%   - Unified Web Service Architecture: FastAPI powering machine JSON REST endpoints and Jinja2/HTMX operator UI.
%   - Route Endpoints:
%       * /predict: multipart image upload, energy OOD filter check, conformal prediction set generation.
%       * /health/live and /health/ready: Kubernetes liveness and readiness probe targets.
%       * /metrics: Prometheus metrics emission.
%       * /train: remote training dispatch endpoint for headless CI/CD execution.
%       * /explain: visual CAM heatmap generation.
%   - HTMX Server-Rendered Interface:
%       * Lightweight HTML fragments swapped via hx-post and hx-target, zero client-side React/Node bloat.
%       * Real-time triage inspection, image drag-and-drop, interactive conformal threshold sliders.
% -------------------------------------------------------------

The TaxonVision-MLOps platform necessitates a robust, dual-purpose service layer capable of fulfilling two distinct operational requirements. First, it must act as a high-throughput machine-to-machine interface for automated ecological sensor networks, exposing strictly validated OpenAPI Representational State Transfer (REST) endpoints. Concurrently, it must provide a responsive, intuitive visual dashboard for human field biologists, taxonomists, and MLOps engineers performing active triage, conformal calibration, and inference evaluation.

This chapter dissects the software architecture bridging these two domains. By leveraging the asynchronous capabilities of FastAPI, the unified service application successfully coalesces machine JSON endpoints with server-rendered Jinja2 HTML templates, augmented by the Hypertext Markup Language extension (HTMX) to eliminate complex client-side state management frameworks.

\section{Unified Web Service Architecture}
\label{sec:unified_web_service_architecture}

The core tenet underlying the TaxonVision application service layer is the `Dual-Personality Web Service Architecture'. Rather than partitioning the system into isolated microservices (e.g., separating the backend inference engine from a decoupled frontend User Interface layer), the application instantiates a single ASGI (Asynchronous Server Gateway Interface) composition root. This unification strategy, implemented within \texttt{src/taxon\_vision/service/api.py}, yields substantial architectural advantages in an MLOps context.

Foremost, it permits shared memory access to heavily vectorised components. In a decoupled architecture, both the machine REST API and the frontend rendering service might redundantly instantiate memory-intensive PyTorch neural networks, ONNX Runtime execution sessions, and large-scale DarwinCore taxonomic catalogues. By unifying the lifespan context within a singular FastAPI instance, these resources are loaded precisely once at application startup into \texttt{app.state}. This shared context drastically reduces the aggregate memory footprint and GPU vRAM consumption on the \texttt{hive-mind} bare-metal Kubernetes nodes, while mitigating initialisation latency across varied endpoints.

The application factory pattern mounts static CSS and JavaScript styling assets (located in \texttt{static/}) whilst orchestrating modular route inclusion. The modularisation is stringently partitioned into functional domains: \texttt{health.py} for orchestration liveness and metrics, \texttt{predict.py} for core inference capabilities, \texttt{train.py} to expose remote continuous training hooks, \texttt{explain.py} for visual attribution maps, and \texttt{views.py} to render the human operator dashboards.

This unified topology also significantly simplifies the Kubernetes pod deployment. A single deployment manifest and corresponding service definition are required. Cloudflare Zero Trust ingress tunnels route external traffic indiscriminately to this service, and the service internally differentiates between API calls demanding JSON serialisation and browser navigation requests expecting HTML payloads.


\subsection{Real-time Event Streaming and Epistemic Triage}

A fundamental requirement for an MLOps platform operating in active learning mode is the seamless identification and routing of highly uncertain or out-of-distribution (OOD) observations to human taxonomists. The TaxonVision architecture facilitates this via the \texttt{/triage} dashboard routes and Server-Sent Events (SSE) integration.

When a field observation triggers an OOD flag (due to elevated Helmholtz free energy) or falls below conformal calibration thresholds, it is automatically enqueued for epistemic triage. The operator dashboard incorporates a real-time listening stream that subscribes to these triage events. Rather than relying on computationally wasteful client-side polling loops, the FastAPI application pushes HTML fragments directly to the connected HTMX clients as soon as a new anomaly is detected.

This hypermedia-driven push mechanism ensures that taxonomists are immediately presented with ambiguous images, CAM attribution heatmaps (retrieved via the \texttt{/explain} endpoint), and the bounding conformal prediction set. The operator can then confirm the novel species, manually adjust the identification, and commit the observation to the active learning pipeline, all without ever leaving the unified server-rendered environment or executing custom client-side JavaScript.

\section{Resilience, Concurrency, and Middleware Engine}
\label{sec:resilience_concurrency_and_middleware}

Serving machine learning visual inference within a highly concurrent web server presents significant computational challenges. The underlying tensor operations (image decoding, normalisation, mathematical matrix multiplication within the ONNX execution session) are fundamentally synchronous and CPU/GPU-bound. Executing these operations directly upon the primary \texttt{asyncio} event loop thread would block all concurrent I/O operations, leading to catastrophic event loop starvation. Under such conditions, instantaneous health probes would timeout, triggering false-positive pod restarts by the Kubernetes kubelet.

To prevent starvation, the TaxonVision architecture implements an asynchronous worker model offloading mechanism. In the inference endpoints, synchronous heavy computation (encapsulated within \texttt{run\_prediction}) is explicitly delegated to a dedicated threadpool executor provided by the AnyIO asynchronous framework via \texttt{anyio.to\_thread.run\_sync}. This preserves the non-blocking nature of the primary event loop, ensuring rapid response times for concurrent requests.

The application architecture embeds a comprehensive middleware engine to handle request tracing and security. Every incoming HTTP request is assigned a unique correlation ID via middleware, propagating through the logging subsystem to ensure that execution crashes, fallback events, and latency spikes are fully traceable back to the originating client payload. Cross-Origin Resource Sharing (CORS) configurations and strict security headers are additionally injected at the middleware boundary to prevent Cross-Site Scripting (XSS) and unauthorised domain access.

\section{API Schemas, Endpoints, and Inference Routing}
\label{sec:api_schemas_and_endpoints}

The machine-to-machine interface of the application is governed by rigid Pydantic v2 validation schemas. These schemas establish an unambiguous data contract between the automated field sensors and the inference engine.

\subsection{Data Contracts and Pydantic Schemas}

The foundational schemas defined in \texttt{schemas.py} encapsulate both taxonomic logic and operational metadata. The \texttt{TaxonPrediction} model enforces strict representation of biological classification, combining the integer DarwinCore taxon ID, the binomial scientific name, the vernacular common name, and the continuous floating-point probability confidence score rounded to four decimal places.

Building upon this, the \texttt{PredictionResponse} schema dictates the overarching response payload for inference requests. Crucially, in accordance with the epistemic safety guidelines detailed in preceding chapters, this schema never permits the return of uncalibrated point estimates in isolation. The payload forces the inclusion of ranked secondary candidates, the scientifically rigorous distribution-free conformal prediction set (nominally providing 95\% marginal coverage), and critical triage boolean flags (\texttt{is\_conformal\_ambiguous}, \texttt{is\_ood\_flagged}, \texttt{requires\_human\_review}). Furthermore, the schema serialises the end-to-end execution latency (\texttt{latency\_ms}), enabling client-side latency monitoring and Pareto boundary verification.

The training pipeline is similarly governed by API contracts. The \texttt{TrainingRequest} schema provides operational parameters to dictate background model head training, including the foundation backbone extractor (e.g., \texttt{mobilenetv4\_conv\_small}), the epoch budget, batch size, and learning rate. The corresponding \texttt{TrainingResponse} structures the convergence outcome, validation accuracy, and final model artifact path.

\subsection{The Inference Predict Endpoint}

The core predictive capability is exposed via the \texttt{/predict} route, constructed within \texttt{predict.py}. This endpoint serves as the primary ingress vector for all high-throughput automated classification tasks originating from the distributed sensor networks. This endpoint accepts multipart \texttt{UploadFile} requests containing the biological imagery.

To ensure stability, the handler enacts a rigorous error-handling and fallback protocol. The raw bytes are read asynchronously before being dispatched to the AnyIO threadpool worker. If inference is successful, the raw prediction dictionaries are transformed by the \texttt{\_build\_live\_prediction} transformer into the validated \texttt{PredictionResponse} model. This transformation logic correlates the raw inference string outputs against the canonical \texttt{taxa\_catalog.yaml} reference to enrich the response with structured taxonomic identifiers.

In the event of an image decoding failure (e.g., due to corrupt binary payloads, raised as a \texttt{UnidentifiedImageError}), the system must distinguish between client errors and application crashes. A graceful fallback mechanism intercepts these specific exceptions, increments a fallback Prometheus telemetry counter, logs a warning, and returns a deterministic, hardcoded fallback prediction (Monarch Butterfly, \textit{Danaus plexippus}). This ensures integration tests and lightweight probe clients can verify endpoint routing, schema validation, and header integrity without possessing binary asset fixtures. Unexpected runtime exceptions (e.g., out-of-memory errors) are deliberately allowed to surface as HTTP 500 errors to alert observability stacks.

\begin{listing}[htpb]
\begin{minted}{python}
@router.post("/predict", response_model=PredictionResponse)
async def predict_species(file: Annotated[UploadFile, File(...)]) -> PredictionResponse:
    t0 = time.perf_counter()
    content = await file.read()
    try:
        pred = await anyio.to_thread.run_sync(run_prediction, content)
        latency = (time.perf_counter() - t0) * 1000.0
        PREDICTION_COUNTER.labels(status="success").inc()
        return _build_live_prediction(pred, latency)
    except (UnidentifiedImageError, OSError, ValueError) as img_err:
        # Handled fallback for corrupt or test payloads
        PREDICTION_COUNTER.labels(status="fallback").inc()
        return _build_fallback_prediction(latency)
\end{minted}
\caption{Simplified excerpt of the offloaded asynchronous prediction handler.}
\label{lst:predict_endpoint}
\end{listing}

\subsection{Kubernetes Liveness, Readiness, and Prometheus Telemetry}

The orchestration fabric demands precise health probing endpoints. The application exposes \texttt{/health} and \texttt{/health/ready} within \texttt{health.py}. These endpoints are intentionally designed to be lightweight, zero-side-effect functions that execute instantaneously. Executing I/O or model inference on these probes would risk false-positive pod terminations under high load. They return a simple JSON status dictating operational liveness.

Simultaneously, the application provides a \texttt{/metrics} endpoint. This route integrates directly with the Prometheus client registry. It bypasses JSON serialisation, instead using \texttt{generate\_latest()} to serialise all instantiated telemetry counters (such as prediction success/fallback rates), histograms (inference latency distribution), and gauges into the standard Prometheus textual exposition format. This pull-based scraping mechanism provides decoupled observability for Grafana dashboards without injecting external network dependencies into the core inference path.

\section{Server-Rendered HTMX Operator Interface}
\label{sec:htmx_operator_interface}

To facilitate human verification, taxonomic review, and active learning curation, the TaxonVision platform exposes a comprehensive operator dashboard. This interface eschews the conventional Single-Page Application (SPA) architecture, instead embracing a Hypermedia-Driven, server-rendered paradigm leveraging Jinja2 templates and HTMX.

\subsection{Rationale for Hypermedia-Driven Architecture}

The selection of a server-rendered HTMX architecture was predicated upon strict operational constraints. In a high-security MLOps environment deployed upon remote edge clusters or secure cloud infrastructure, minimising the software supply chain attack surface is paramount. The traditional SPA ecosystem (e.g., React, Vue, Angular) necessitates complex JavaScript build toolchains, mandating Node.js runtimes, \texttt{npm} package managers, Webpack bundlers, and thousands of transient dependencies.

By actively banning the Node.js runtime and \texttt{npm} ecosystem from the production container build, the TaxonVision architecture drastically reduces its vulnerability profile. The hypermedia architecture shifts state management and DOM (Document Object Model) generation entirely to the Python backend. The client requires only a minuscule, zero-dependency HTMX JavaScript library served as a static asset. This ensures rapid, instant server-rendered markup delivery, drastically simplifying the frontend complexity while maintaining a highly interactive user experience.

\subsection{Dynamic DOM Swapping and Route Implementation}

The HTMX library augments standard HTML elements with declarative attributes (\texttt{hx-post}, \texttt{hx-target}, \texttt{hx-swap}), enabling asynchronous server requests and partial DOM updates without full page reloads.

The primary dashboard is rendered by the \texttt{/} route within \texttt{views.py}. This endpoint binds the FastAPI request context to the Jinja2 \texttt{index.html} template. The template constructs a responsive user interface with an image upload form explicitly configured for HTMX interaction.

When a user submits an observation photograph, the form does not execute a standard POST request that would refresh the entire browser window. Instead, \texttt{hx-post="/ui/predict-htmx"} intercepts the submission. The image binary is asynchronously transmitted to the \texttt{predict\_htmx} endpoint.

\begin{listing}[htpb]
\begin{minted}{html}
  <form id="upload-form"
        method="post"
        action="/ui/predict-htmx"
        hx-post="/ui/predict-htmx"
        hx-target="#prediction-result"
        hx-swap="innerHTML"
        hx-indicator="#loading-indicator"
        enctype="multipart/form-data">
\end{minted}
\caption{HTMX declarative attributes binding the form submission to partial DOM swaps.}
\label{lst:htmx_form}
\end{listing}

The \texttt{predict\_htmx} route functions similarly to the JSON machine endpoint, offloading inference to the AnyIO threadpool. Crucially, however, instead of returning a Pydantic JSON string, it constructs a Jinja2 rendering context containing the taxonomic classifications, conformal sets, and triage flags. It then evaluates the \texttt{partials/prediction\_card.html} template.

This partial HTML fragment represents solely the result card (containing scientific nomenclature badges, confidence percentages, and explicitly rendered Out-Of-Distribution alert warnings). The FastAPI server transmits this small HTML snippet back to the client browser. HTMX automatically detects the response, identifies the \texttt{hx-target="\#prediction-result"} directive, and seamlessly swaps the inner HTML of the target element with the newly generated prediction card.

This seamless partial swapping paradigm provides an interactive, SPA-like user experience characterised by instantaneous visual feedback and fluid transitions. It achieves this (including loading indicators managed via \texttt{hx-indicator}) while retaining the architectural purity and security profile of a traditional server-rendered application. Real-time event streaming and Server-Sent Events (SSE) can be trivially integrated into this architecture, pushing newly referred observations from the active learning queue directly to the operator dashboard without requiring continuous client-side polling loops.
